\chapter{Introduction}
\label{sec:intro}
\section{Purpose}
High-performance computing dataset sizes are
increasing.\cite{reed_exascale_2015}. As \citeauthor{underwood_fraz_2020} state
in \cite{underwood_fraz_2020}, ``Today’s scientific research applications
produce volumes of data too large to be stored, transferred, and analyzed
efficiently because of limited storage space and potential bottlenecks in I/O
systems.'' \cite{foster_computing_2017} provides three examples of applications
generating large volumes of data:
\begin{itemize}
\item Climate science simulations checkpointing generating \qty{260}{\tera\byte}
  every 16~seconds.
\item Singe XGC (``X-point included Gyrokinetics Code'') ITER fusion models
  producing ``hundreds of petabytes''\cite{foster_computing_2017}.
\item Femtosecond-level material science simulations producing ``exabytes on
  exascale computers''\cite{foster_computing_2017}.
\end{itemize}

\cite{foster_computing_2017} notes this as a problem due to 2024-projected (at
the time) compute rate outperforming disk write rate by 50 times. Because I/O
throughput scales more slowly than computational speeds, ``We can no longer
output every piece of information that we might ever possibly
want''\cite{foster_computing_2017}.

\section{OpenZL}
The meat and potatoes of this thesis rests on OpenZL and its innovations. What
is OpenZL? OpenZL is a graph-based compressor released by \metainc{} in a
whitepaper, ``OpenZL: A Graph-Based Model for
Compression''\cite{collet_openzl_2025}, in 2025. It allows for greater
performance over generic compressors without the complexity and challenges of
application-specific compressors. The compressed data format is ``a
self-describing wire format, \ldots which can be decompressed by a universal
decoder''\cite{collet_openzl_2025}. Because of this, it allows compressors
tailored to datasets with universal decoding. The whitepaper lists the following
benefits:
\begin{itemize}
\item Minimal upfront investment
\item Flexibility of input data formats
\item Minimized security surface
\item Easy iteration
\item Fast deployment
\item Support for high-cardinality data
\end{itemize}
